% Chapter 3 Introduction (untitled opening paragraphs)
% Funnel structure: broad context → specific problem → chapter roadmap

The gamma-ray sky described in the previous chapter presents a formidable data analysis challenge.
Faint signals from potential dark matter annihilation or decay are embedded in bright, spatially structured astrophysical backgrounds; source populations overlap in both position and spectral shape; and instrumental effects---finite angular resolution, energy-dependent sensitivity, detection thresholds---leave a large fraction of the emitting population unresolved or unclassified.
Extracting new physics from this environment requires moving beyond standard signal-over-background methods toward a more comprehensive statistical toolkit.

The analyses presented in this thesis draw on techniques spanning classical statistics, modern machine learning, and domain-specific tools developed for astrophysical inference.
Each method addresses a different aspect of the extraction problem.
Frequentist and Bayesian approaches provide the language for parameter estimation and hypothesis testing.
Simulation-based inference bypasses the intractable likelihoods that arise in complex forward models, while neural networks learn directly from high-dimensional data when explicit physical models are unavailable.
Domain adaptation corrects for the systematic mismatches between simulated training data and real observations.
Finally, angular cross-correlations exploit the connection between gamma-ray emission and the large-scale distribution of matter, accessing collective signals invisible to single-instrument analyses.

This chapter introduces the mathematical framework used throughout the thesis.
Each method is presented at a conceptual level, with \blue{forward references to the chapters} where the full technical details appear.
We begin with the foundations of frequentist and Bayesian inference in Section~\ref{sec:3.1}, establishing the likelihood function, maximum likelihood estimation, Bayesian posteriors, and hypothesis testing.
Section~\ref{sec:3.2} introduces the simulation-based inference paradigm for problems where the likelihood cannot be evaluated analytically.
% Section~\ref{sec:3.3} covers the machine learning tools---convolutional neural networks, density estimation, and mixture models---used for regression and classification tasks on gamma-ray data.
\blue{Section~\ref{sec:3.3} covers the machine learning tools---convolutional neural networks and density estimation---used for regression and classification tasks on gamma-ray data.}
Section~\ref{sec:3.4} formalizes the dataset shift problem that arises when training and target distributions differ.
Finally, Section~\ref{sec:cross_correlations} develops the angular cross-correlation formalism that connects gamma-ray observations to gravitational tracers of dark matter structure at cosmological scales.
